\chapter{Lazy Iterators}\label{ch:lazy-iterators}

\index{iterators}\index{lazy evaluation}\index{Iterator type}

Imagine you have a log file with ten million lines, and you need the first twenty that contain the word \texttt{"ERROR"}.
One approach: read all ten million lines into an array, filter it, then slice off the first twenty.
That works, but it loads the entire file into memory and does far more work than necessary.

There is a better way.
\emph{Iterators} let you process data one element at a time, on demand, without ever building the full collection in memory.
In Lattice, iterators are \emph{lazy}---they produce values only when asked.
This chapter shows you how to create iterators, chain transformations, consume results, and build your own custom iterators from scratch.

%% ───────────────────────────────────────────────────────────
\section{The Iterator Protocol}\label{sec:iterator-protocol}
\index{iterators!protocol}\index{iter()}\index{iter\_next()}

At its heart, an iterator is anything that can produce a sequence of values, one at a time, until it is exhausted.
Lattice's built-in \lstinline|iter()| function converts a collection into an \lstinline|Iterator| value, and the \lstinline|.next()| method advances it by one step.

\begin{lstlisting}[language=Lattice, caption={The basic iterator protocol}]
let colors = ["crimson", "teal", "gold"]
let it = iter(colors)

print(it.next())  // crimson
print(it.next())  // teal
print(it.next())  // gold
print(it.next())  // nil
\end{lstlisting}

Each call to \lstinline|.next()| returns the next value.
When the iterator is exhausted, it returns \lstinline|nil|.

\begin{defnbox}{Iterator}
An \emph{iterator} is a value of type \lstinline|Iterator| that encapsulates:
\begin{enumerate}[nosep]
  \item Some internal \textbf{state} (a position, a counter, a reference to a source).
  \item A \textbf{next function} that, given the current state, either produces the next value or signals that the sequence is done.
\end{enumerate}
Once an iterator is exhausted, every subsequent call to \lstinline|.next()| returns \lstinline|nil|.
Iterators are single-pass: you cannot rewind them.
\end{defnbox}

\subsection{What Can You Iterate Over?}\label{subsec:iterable-types}
\index{iter()!supported types}

The \lstinline|iter()| function accepts arrays, maps, sets, strings, ranges, and even existing iterators (which it returns unchanged):

\begin{lstlisting}[language=Lattice, caption={Creating iterators from different collections}]
// Array: yields elements in order
let arr_it = iter([10, 20, 30])

// String: yields individual characters
let str_it = iter("hello")
print(str_it.next())  // h
print(str_it.next())  // e

// Range: yields integers lazily
let range_it = iter(0..5)
print(range_it.next())  // 0
print(range_it.next())  // 1

// Map: yields keys
let scores = Map::new()
scores.set("alice", 95)
scores.set("bob", 87)
let map_it = iter(scores)
print(map_it.next())  // alice (or bob -- map order is not guaranteed)
\end{lstlisting}

\begin{notebox}{Iterators Are Single-Pass}
Once you call \lstinline|.next()| and consume a value, it is gone.
If you need to iterate over the same data twice, create a new iterator with \lstinline|iter()| again---or \lstinline|.collect()| the results into an array first.
\end{notebox}

\subsection{Using Iterators in For Loops}\label{subsec:iter-for-loops}
\index{iterators!for loops}

You rarely call \lstinline|.next()| by hand.
Lattice's \lstinline|for| loop consumes any iterable automatically:

\begin{lstlisting}[language=Lattice, caption={Iterators and for loops}]
let temperatures = [72.1, 68.4, 75.0, 71.3]

// The for loop calls iter() and .next() behind the scenes
for temp in temperatures {
    print("Reading: " + to_string(temp))
}
\end{lstlisting}

When you write \lstinline|for x in collection|, Lattice creates an iterator from the collection and repeatedly calls \lstinline|.next()| until it returns \lstinline|nil|.
This is the same protocol you just learned---the \lstinline|for| loop is syntactic sugar for the manual \lstinline|.next()| loop.

%% ───────────────────────────────────────────────────────────
\section{Sources: Where Iterators Come From}\label{sec:iterator-sources}
\index{iterators!sources}

The \lstinline|iter()| function wraps existing collections, but Lattice also provides dedicated \emph{source} constructors that create iterators from scratch.

\subsection{iter() --- From Collections}\label{subsec:iter-from-collection}

As we saw above, \lstinline|iter()| is the universal entry point.
It accepts arrays, maps, sets, strings, ranges, and iterators.
Under the hood, \lstinline|iter()| dispatches to specialized constructors like \lstinline|iter_from_array|, \lstinline|iter_from_map|, and \lstinline|iter_from_range| in the C runtime (see \filename{src/iterator.c}).

\subsection{range\_iter() --- Lazy Integer Sequences}\label{subsec:range-iter}
\index{range\_iter()}\index{iterators!range}

While \lstinline|iter(0..1000)| converts a range to an iterator, the dedicated \lstinline|range_iter()| function gives you more control---including a custom step:

\begin{lstlisting}[language=Lattice, caption={Creating range iterators}]
// Count from 0 to 9
let digits = range_iter(0, 10)
print(digits.collect())  // [0, 1, 2, 3, 4, 5, 6, 7, 8, 9]

// Even numbers from 0 to 20
let evens = range_iter(0, 20, 2)
print(evens.collect())  // [0, 2, 4, 6, 8, 10, 12, 14, 16, 18]

// Countdown
let countdown = range_iter(5, 0, -1)
print(countdown.collect())  // [5, 4, 3, 2, 1]
\end{lstlisting}

The third argument is the step.
If omitted, it defaults to \lstinline|1| for ascending ranges or \lstinline|-1| for descending ranges.
No array is allocated---the iterator computes each value on the fly.

\begin{tipbox}{When to Use range\_iter() vs.\ a Range Literal}
If you need a custom step or you are building a pipeline of iterator transformations, use \lstinline|range_iter()|.
For straightforward \lstinline|for| loops, the range literal (\lstinline|0..10|) is more concise and equally lazy.
\end{tipbox}

\subsection{repeat\_iter() --- The Same Value, Over and Over}\label{subsec:repeat-iter}
\index{repeat\_iter()}\index{iterators!repeat}

Need a sequence of identical values?
\lstinline|repeat_iter()| yields the same value a specified number of times---or infinitely if you omit the count:

\begin{lstlisting}[language=Lattice, caption={Repeating a value}]
// Repeat a value 4 times
let zeros = repeat_iter(0, 4)
print(zeros.collect())  // [0, 0, 0, 0]

// Infinite repetition (be careful!)
let ones = repeat_iter(1)
// ones.collect() would run forever --- use .take() instead
print(ones.take(5).collect())  // [1, 1, 1, 1, 1]
\end{lstlisting}

\begin{warnbox}{Infinite Iterators}
Calling \lstinline|.collect()| on an infinite iterator will consume all available memory and hang your program.
Always pair infinite sources with a limiting transformer like \lstinline|.take()| before consuming.
\end{warnbox}

%% ───────────────────────────────────────────────────────────
\section{Transformers: Reshaping the Stream}\label{sec:iterator-transformers}
\index{iterators!transformers}

The real power of iterators emerges when you chain \emph{transformers}---methods that wrap one iterator in another, building up a processing pipeline without doing any work until you ask for results.

\subsection{.map() --- Transform Each Element}\label{subsec:iter-map}
\index{iterators!map}

The \lstinline|.map()| method applies a closure to every element the iterator produces:

\begin{lstlisting}[language=Lattice, caption={Mapping over an iterator}]
let prices = [19.99, 45.50, 12.00, 89.95]
let with_tax = iter(prices)
    .map(|price| { price * 1.08 })
    .collect()

print(with_tax)
// [21.5892, 49.14, 12.96, 97.146]
\end{lstlisting}

The closure receives each value and returns its transformation.
Because \lstinline|.map()| returns a new iterator, no work happens until \lstinline|.collect()| pulls values through.

\subsection{.filter() --- Keep Only What Matters}\label{subsec:iter-filter}
\index{iterators!filter}

The \lstinline|.filter()| method takes a predicate closure and yields only the elements for which it returns \lstinline|true|:

\begin{lstlisting}[language=Lattice, caption={Filtering an iterator}]
let readings = [3.1, 7.8, 2.4, 9.0, 5.5, 1.2]
let high_readings = iter(readings)
    .filter(|r| { r > 5.0 })
    .collect()

print(high_readings)  // [7.8, 9.0, 5.5]
\end{lstlisting}

Internally, the filter iterator loops through the inner source, skipping elements that don't match, until it finds one that does---or the source is exhausted.
You can see this logic in \filename{src/iterator.c} where \lstinline|iter_filter_next| calls the inner iterator in a loop.

\subsection{.take() --- Limit the Count}\label{subsec:iter-take}
\index{iterators!take}

The \lstinline|.take(n)| transformer yields at most \lstinline|n| elements, then stops:

\begin{lstlisting}[language=Lattice, caption={Taking the first few elements}]
let first_three = range_iter(0, 1000000)
    .take(3)
    .collect()

print(first_three)  // [0, 1, 2]
\end{lstlisting}

Despite the range spanning a million integers, only three are ever produced.
This is laziness in action: the remaining 999{,}997 values are never computed.

\subsection{.skip() --- Jump Ahead}\label{subsec:iter-skip}
\index{iterators!skip}

The companion to \lstinline|.take()|, the \lstinline|.skip(n)| transformer discards the first \lstinline|n| elements and yields the rest:

\begin{lstlisting}[language=Lattice, caption={Skipping elements}]
let after_header = iter(["name", "age", "city", "Alice", "30", "Portland"])
    .skip(3)
    .collect()

print(after_header)  // [Alice, 30, Portland]
\end{lstlisting}

You can combine \lstinline|.skip()| and \lstinline|.take()| to extract a window from a sequence:

\begin{lstlisting}[language=Lattice, caption={Pagination with skip and take}]
let all_items = range_iter(1, 101)   // items 1 through 100
let page_size = 10
let page_number = 3

let page = all_items
    .skip((page_number - 1) * page_size)
    .take(page_size)
    .collect()

print(page)  // [21, 22, 23, 24, 25, 26, 27, 28, 29, 30]
\end{lstlisting}

\subsection{.enumerate() --- Track the Index}\label{subsec:iter-enumerate}
\index{iterators!enumerate}

When you need both the position and the value, \lstinline|.enumerate()| wraps each element in a two-element array \lstinline|[index, value]|:

\begin{lstlisting}[language=Lattice, caption={Enumerating an iterator}]
let languages = ["Lattice", "Rust", "Python", "Go"]
let indexed = iter(languages).enumerate()

for pair in indexed.collect() {
    print(to_string(pair[0]) + ": " + pair[1])
}
// 0: Lattice
// 1: Rust
// 2: Python
// 3: Go
\end{lstlisting}

The index always starts at zero.

\subsection{.zip() --- Walk Two Iterators in Lockstep}\label{subsec:iter-zip}
\index{iterators!zip}

The \lstinline|.zip()| transformer pairs elements from two iterators into \lstinline|[left, right]| arrays.
It stops as soon as either iterator is exhausted:

\begin{lstlisting}[language=Lattice, caption={Zipping two iterators}]
let names = iter(["Alice", "Bob", "Carol"])
let scores = iter([95, 87, 92])

let paired = names.zip(scores).collect()
print(paired)
// [[Alice, 95], [Bob, 87], [Carol, 92]]
\end{lstlisting}

Zip is especially useful for combining parallel data sources:

\begin{lstlisting}[language=Lattice, caption={Merging labels with values}]
let labels = iter(["Temperature", "Humidity", "Pressure"])
let values = iter([72.1, 45.0, 1013.25])
let units = iter(["F", "%", "hPa"])

// Zip labels with values, then map to formatted strings
let report = labels.zip(values).zip(units)
    .map(|entry| {
        let label_value = entry[0]
        let unit = entry[1]
        label_value[0] + ": " + to_string(label_value[1]) + " " + unit
    })
    .collect()

for line in report {
    print(line)
}
// Temperature: 72.1 F
// Humidity: 45.0 %
// Pressure: 1013.25 hPa
\end{lstlisting}

\subsection{Chaining Transformers}\label{subsec:chaining-transformers}
\index{iterators!chaining}

Transformers return iterators, so you can chain them into expressive pipelines:

\begin{lstlisting}[language=Lattice, caption={A multi-step pipeline}]
let words = ["Lattice", "is", "a", "language", "for", "the", "curious"]

let result = iter(words)
    .filter(|w| { len(w) > 2 })
    .map(|w| { w + "!" })
    .enumerate()
    .take(3)
    .collect()

print(result)
// [[0, Lattice!], [1, language!], [2, for!]]
\end{lstlisting}

No intermediate arrays are created.
Each value flows through the entire chain---from source to \lstinline|.filter()| to \lstinline|.map()| to \lstinline|.enumerate()| to \lstinline|.take()|---one at a time, on demand.

\begin{tipbox}{Reading Pipelines}
Read iterator chains from top to bottom, like a recipe:
\begin{enumerate}[nosep]
  \item Start with the words.
  \item Keep only words longer than two characters.
  \item Append an exclamation mark.
  \item Number them.
  \item Take the first three.
  \item Collect into an array.
\end{enumerate}
\end{tipbox}

%% ───────────────────────────────────────────────────────────
\section{Consumers: Triggering the Work}\label{sec:iterator-consumers}
\index{iterators!consumers}

Transformers are lazy---they build up a chain of work but produce no output on their own.
\emph{Consumers} are the methods that actually pull values through the pipeline and produce a final result.

\subsection{.collect() --- Materialize into an Array}\label{subsec:iter-collect}
\index{iterators!collect}

We have been using \lstinline|.collect()| throughout this chapter.
It drains the iterator and gathers every element into an array:

\begin{lstlisting}[language=Lattice, caption={Collecting an iterator}]
let squares = range_iter(1, 6)
    .map(|n| { n * n })
    .collect()

print(squares)  // [1, 4, 9, 16, 25]
\end{lstlisting}

Under the hood, \lstinline|iter_collect| in \filename{src/iterator.c} starts with a small buffer (capacity 8) and doubles it as needed, appending each value from \lstinline|iter_next| until the iterator signals done.
The alias \lstinline|.to_array()| does exactly the same thing.

\subsection{.reduce() --- Fold into a Single Value}\label{subsec:iter-reduce}
\index{iterators!reduce}\index{iterators!fold}

The \lstinline|.reduce()| method collapses an iterator into a single accumulated value.
It takes a closure and an initial accumulator:

\begin{lstlisting}[language=Lattice, caption={Reducing an iterator}]
let total = iter([10, 20, 30, 40])
    .reduce(|acc, val| { acc + val }, 0)

print(total)  // 100
\end{lstlisting}

The closure receives two arguments: the running accumulator and the current element.
The second argument to \lstinline|.reduce()| is the initial value for the accumulator.

Here is a more involved example---computing the maximum value:

\begin{lstlisting}[language=Lattice, caption={Finding the maximum with reduce}]
let temperatures = [68.2, 71.5, 65.0, 73.8, 70.1]

let hottest = iter(temperatures)
    .reduce(|max_temp, temp| {
        if temp > max_temp { temp } else { max_temp }
    }, temperatures[0])

print(hottest)  // 73.8
\end{lstlisting}

\subsection{.any() and .all() --- Short-Circuit Boolean Tests}\label{subsec:iter-any-all}
\index{iterators!any}\index{iterators!all}

These two consumers answer yes-or-no questions about the sequence.
Both \emph{short-circuit}: they stop consuming as soon as the answer is determined.

\begin{lstlisting}[language=Lattice, caption={Testing with .any() and .all()}]
let grades = [88, 92, 75, 95, 61]

let has_failing = iter(grades).any(|g| { g < 65 })
print(has_failing)  // true (stops after finding 61)

let all_passing = iter(grades).all(|g| { g >= 60 })
print(all_passing)  // true
\end{lstlisting}

The \lstinline|.any()| consumer returns \lstinline|true| as soon as it finds a matching element, without examining the rest.
The \lstinline|.all()| consumer returns \lstinline|false| as soon as it finds a non-matching element.
If the iterator is empty, \lstinline|.any()| returns \lstinline|false| and \lstinline|.all()| returns \lstinline|true|---consistent with the mathematical conventions for existential and universal quantification.

\subsection{.count() --- How Many?}\label{subsec:iter-count}
\index{iterators!count}

The \lstinline|.count()| consumer drains the iterator and returns the number of elements:

\begin{lstlisting}[language=Lattice, caption={Counting elements}]
let long_words = iter(["cat", "elephant", "dog", "rhinoceros", "ox"])
    .filter(|w| { len(w) > 4 })
    .count()

print(long_words)  // 2
\end{lstlisting}

\begin{notebox}{Consumers Are Destructive}
Calling a consumer like \lstinline|.collect()|, \lstinline|.reduce()|, or \lstinline|.count()| drains the iterator.
If you try to use the same iterator afterwards, it will produce no more values.
To reuse, build a fresh iterator or collect into an array first.
\end{notebox}

\subsection{Summary of Consumers}\label{subsec:consumer-summary}

\begin{table}[htbp]
\centering
\caption{Iterator consumer methods}\label{tab:iter-consumers}
\begin{tabular}{@{}lll@{}}
\toprule
\textbf{Method} & \textbf{Returns} & \textbf{Description} \\
\midrule
\lstinline|.collect()| & \lstinline|Array| & Gather all elements into an array \\
\lstinline|.to_array()| & \lstinline|Array| & Alias for \lstinline|.collect()| \\
\lstinline|.reduce(fn, init)| & \lstinline|Any| & Fold with accumulator \\
\lstinline|.any(predicate)| & \lstinline|Bool| & True if any element matches \\
\lstinline|.all(predicate)| & \lstinline|Bool| & True if all elements match \\
\lstinline|.count()| & \lstinline|Int| & Count elements (consumes all) \\
\bottomrule
\end{tabular}
\end{table}

%% ───────────────────────────────────────────────────────────
\section{Building Your Own Iterators}\label{sec:custom-iterators}
\index{iterators!custom}\index{fn standard library}

The built-in iterator methods cover the common cases, but sometimes you need a custom sequence---Fibonacci numbers, lines from a database cursor, or values generated by a complex algorithm.

Lattice's \filename{lib/fn.lat} standard library module provides a protocol for building custom lazy sequences entirely in Lattice code.
Import it, and you get access to a rich set of sequence constructors and transformers.

\subsection{The Sequence Protocol}\label{subsec:seq-protocol}
\index{lazy sequences}\index{fn standard library!sequences}

A lazy sequence in the \lstinline|fn| module is a map with two keys:
\begin{itemize}[nosep]
  \item \lstinline|"state"|: the current position or state of the iterator (any value).
  \item \lstinline|"next"|: a function that takes the current state and returns a step map.
\end{itemize}

Each \emph{step} is itself a map:
\begin{itemize}[nosep]
  \item \lstinline|\{"value": v, "done": false, "state": next_state\}| --- yields value \lstinline|v| and advances.
  \item \lstinline|\{"done": true\}| --- the sequence is exhausted.
\end{itemize}

\begin{defnbox}{Lazy Sequence (fn module)}
A lazy sequence threads state explicitly through each step.
This design is compatible with Lattice's value semantics: no hidden mutable references, no shared state.
The next function is pure---given the same state, it always produces the same step.
\end{defnbox}

\subsection{Building a Fibonacci Iterator}\label{subsec:fibonacci-iter}

Let's build a lazy Fibonacci sequence using the \lstinline|fn| module:

\begin{lstlisting}[language=Lattice, caption={A custom Fibonacci iterator}]
import "lib/fn" as F

// State is [current, next_val]
let fibs = F.iterate([0, 1], |pair| {
    [pair[1], pair[0] + pair[1]]
})

// iterate() yields the state at each step,
// so we fmap to extract just the first element
let fib_numbers = F.fmap(fibs, |pair| { pair[0] })

// Take the first 10 Fibonacci numbers
let result = F.collect(F.take(fib_numbers, 10))
print(result)  // [0, 1, 1, 2, 3, 5, 8, 13, 21, 34]
\end{lstlisting}

The key function here is \lstinline|F.iterate(seed, f)|, which produces an infinite sequence: \lstinline|seed|, \lstinline|f(seed)|, \lstinline|f(f(seed))|, and so on.
We start with \lstinline|[0, 1]| and advance by shifting the pair forward.

\subsection{Custom Sequences from Scratch}\label{subsec:custom-from-scratch}

You can also build a sequence manually by providing a state and a next function.
Here is a sequence that yields powers of two:

\begin{lstlisting}[language=Lattice, caption={Building a sequence from scratch}]
import "lib/fn" as F

fn powers_of_two() -> Map {
    // State is the current power
    let initial_state = 1
    let next_fn = |current| {
        let step = Map::new()
        step.set("value", current)
        step.set("done", false)
        step.set("state", current * 2)
        return step
    }

    let seq = Map::new()
    seq.set("state", initial_state)
    seq.set("next", next_fn)
    return seq
}

let result = F.collect(F.take(powers_of_two(), 8))
print(result)  // [1, 2, 4, 8, 16, 32, 64, 128]
\end{lstlisting}

The pattern is always the same: define an initial state, write a function that takes the current state and returns a step map, then combine them into a sequence map.

\subsection{A Practical Example: Paginated API Results}\label{subsec:paginated-api}

Custom iterators shine when the data source is external or unbounded.
Here is a sketch of a lazy sequence that fetches pages from an API:

\begin{lstlisting}[language=Lattice, caption={Lazy pagination (conceptual)}]
import "lib/fn" as F

fn paginated_users(base_url: String) -> Map {
    // State is the current page number
    return F.iterate(1, |page| { page + 1 })
}

// In a real app, you would use fmap to fetch each page:
// let pages = F.fmap(paginated_users("/api/users"), |page| {
//     http_get(base_url + "?page=" + to_string(page))
// })
//
// Then take pages until the response is empty.
\end{lstlisting}

The lazy approach means you never fetch a page you don't need.

\subsection{Composing fn Module Sequences}\label{subsec:composing-fn-sequences}

The \lstinline|fn| module provides the same family of transformers as the built-in iterator methods, but as standalone functions that work with the sequence protocol:

\begin{lstlisting}[language=Lattice, caption={Composing sequences with the fn module}]
import "lib/fn" as F

// Start with a range, filter, transform, and collect
let result = F.collect(
    F.take(
        F.fmap(
            F.select(F.range(1, 100), |x| { x % 3 == 0 }),
            |x| { x * x }
        ),
        5
    )
)

print(result)  // [9, 36, 81, 144, 225]
\end{lstlisting}

This computes the squares of the first five multiples of three.
Every step is lazy---the range never produces 100 values, only as many as needed to fill the first five matches.

\begin{table}[htbp]
\centering
\caption{fn module sequence functions}\label{tab:fn-seq-functions}
\begin{tabular}{@{}lll@{}}
\toprule
\textbf{Function} & \textbf{Kind} & \textbf{Description} \\
\midrule
\lstinline|F.range(start, end, step?)| & Source & Lazy integer range \\
\lstinline|F.from_array(arr)| & Source & Lazy sequence from array \\
\lstinline|F.repeat(value)| & Source & Infinite repetition \\
\lstinline|F.iterate(seed, f)| & Source & Infinite \lstinline|seed, f(seed), ...| \\
\lstinline|F.fmap(seq, f)| & Transformer & Transform each element \\
\lstinline|F.select(seq, pred)| & Transformer & Keep matching elements \\
\lstinline|F.take(seq, n)| & Transformer & First \lstinline|n| elements \\
\lstinline|F.drop(seq, n)| & Transformer & Skip first \lstinline|n| \\
\lstinline|F.take_while(seq, pred)| & Transformer & Yield while predicate holds \\
\lstinline|F.zip(seq1, seq2)| & Transformer & Pair elements from two sequences \\
\lstinline|F.collect(seq)| & Consumer & Collect into array \\
\lstinline|F.fold(seq, init, f)| & Consumer & Reduce with accumulator \\
\lstinline|F.count(seq)| & Consumer & Count elements \\
\bottomrule
\end{tabular}
\end{table}

\begin{notebox}{Two Iterator Systems}
Lattice has two complementary iteration systems:
\begin{enumerate}[nosep]
  \item \textbf{Built-in iterators} (\lstinline|iter()|, \lstinline|.map()|, \lstinline|.filter()|, etc.)\ --- implemented in C (\filename{src/iterator.c}), accessed via method chaining. Best for everyday use.
  \item \textbf{fn module sequences} (\lstinline|F.range()|, \lstinline|F.fmap()|, \lstinline|F.select()|, etc.)\ --- implemented in pure Lattice (\filename{lib/fn.lat}), using a map-based protocol. Best for custom sequences and functional composition.
\end{enumerate}
Both are lazy. The built-in system is faster (native C); the \lstinline|fn| module is more flexible (user-extensible).
\end{notebox}

%% ───────────────────────────────────────────────────────────
\section{Why Laziness Matters}\label{sec:why-laziness}
\index{lazy evaluation!benefits}\index{iterators!performance}

We have seen iterators in action, but let's make the case for \emph{why} you should prefer them over eager array operations.

\subsection{Memory Efficiency}\label{subsec:lazy-memory}

Consider processing a large dataset:

\begin{lstlisting}[language=Lattice, caption={Eager vs.\ lazy memory usage}]
// Eager: creates three intermediate arrays
let data = range_iter(0, 1000000).collect()  // 1M element array
let filtered = data.filter(|x| { x % 7 == 0 })  // another array
let mapped = filtered.map(|x| { x * x })  // yet another
let result = mapped[0..5]  // finally, what we wanted

// Lazy: creates zero intermediate arrays
let result_lazy = range_iter(0, 1000000)
    .filter(|x| { x % 7 == 0 })
    .map(|x| { x * x })
    .take(5)
    .collect()

print(result_lazy)  // [0, 49, 196, 441, 784]
\end{lstlisting}

The eager version allocates three full arrays before extracting five values.
The lazy version allocates one small array (the final result) and produces only the 35 integers needed to fill it.

\subsection{Computation Efficiency}\label{subsec:lazy-computation}

Laziness also saves computation.
When \lstinline|.take(5)| has collected its five elements, the entire pipeline stops---no more filtering, no more mapping, and the range never produces its remaining elements.

This is especially valuable when the source is expensive (network calls, file I/O) or the filter is highly selective.

\subsection{Infinite Sequences}\label{subsec:infinite-sequences}

Some sequences have no natural end---the natural numbers, a stream of sensor readings, lines from a log file that is still being written.
Lazy iterators handle these gracefully because they never try to materialize the whole sequence:

\begin{lstlisting}[language=Lattice, caption={Working with infinite sequences}]
import "lib/fn" as F

// The natural numbers
let naturals = F.range(1, 9999999)

// First 5 perfect squares
let squares = F.collect(
    F.take(
        F.fmap(naturals, |n| { n * n }),
        5
    )
)
print(squares)  // [1, 4, 9, 16, 25]

// First 5 numbers divisible by both 3 and 7
let div_3_and_7 = F.collect(
    F.take(
        F.select(F.range(1, 9999999), |n| { n % 3 == 0 && n % 7 == 0 }),
        5
    )
)
print(div_3_and_7)  // [21, 42, 63, 84, 105]
\end{lstlisting}

\subsection{When to Stay Eager}\label{subsec:when-eager}

Laziness is not always the right choice.
If you need to access elements by index, sort the collection, or iterate over the data multiple times, collecting into an array first is the pragmatic approach.
Iterators are single-pass and sequential---they trade random access for efficiency.

\begin{tipbox}{A Rule of Thumb}
Use iterators when you are \emph{flowing} data through a pipeline: read, transform, consume.
Use arrays when you need to \emph{hold} data: index, sort, search, iterate multiple times.
\end{tipbox}

%% ───────────────────────────────────────────────────────────
\section{Under the Hood}\label{sec:iterator-internals}
\index{iterators!internals}

If you are curious about how iterators work at the C level, here is a brief tour.

Each built-in iterator is a \lstinline|LatValue| of type \lstinline|VAL_ITERATOR|.
It stores three things:
\begin{itemize}[nosep]
  \item A \textbf{state pointer} (\lstinline|void *state|)---an opaque blob of data specific to the iterator kind (array index, range counter, wrapped inner iterator, etc.).
  \item A \textbf{next function} (\lstinline|next_fn|)---a C function pointer that advances the state and returns the next value.
  \item A \textbf{free function} (\lstinline|free_fn|)---cleans up the state when the iterator is no longer needed.
\end{itemize}

The \lstinline|iter_next| call (defined in \filename{include/iterator.h}) is a single function pointer dispatch:

\begin{quote}
\texttt{return iter->as.iterator.next\_fn(iter->as.iterator.state, \&done);}
\end{quote}

Transformer iterators like \lstinline|iter_take| wrap an inner iterator by storing it in their state struct (e.g., \lstinline|IterTakeState| contains the inner iterator and a remaining count).
When their \lstinline|next_fn| is called, they delegate to the inner iterator's \lstinline|next_fn|, applying their transformation.
This forms a chain of function pointer calls---lightweight and allocation-free during iteration.

\begin{notebox}{Iterator Ownership}
Transformer constructors like \lstinline|iter_take()| and \lstinline|iter_map_transform()| \emph{take ownership} of the inner iterator.
Once you pass an iterator to a transformer, you should not use the original.
The built-in method-chaining API handles this naturally---each \lstinline|.take()| or \lstinline|.map()| call consumes the receiver and returns a new iterator.
\end{notebox}

%% ───────────────────────────────────────────────────────────
\section{Exercises}\label{sec:iter-exercises}

\begin{enumerate}
  \item \textbf{Sum of Squares.}
    Use \lstinline|range_iter|, \lstinline|.map()|, and \lstinline|.reduce()| to compute the sum of squares of the integers from 1 to 100.
    What answer do you get?

  \item \textbf{Custom Collatz Iterator.}
    The Collatz sequence for a starting number \lstinline|n| is: if \lstinline|n| is even, the next value is \lstinline|n / 2|; if odd, it is \lstinline|3 * n + 1|.
    The sequence ends when it reaches 1.
    Using the \lstinline|fn| module's \lstinline|iterate()| and \lstinline|take_while()|, build a lazy Collatz sequence starting from 27 and collect it into an array.
    How many steps does it take?

  \item \textbf{Zip and Reduce.}
    Given two arrays of equal length---one of item names and one of prices---use \lstinline|iter()|, \lstinline|.zip()|, and \lstinline|.reduce()| to compute the total price.
    For example: names \lstinline|["apple", "bread", "milk"]| and prices \lstinline|[1.20, 3.50, 2.80]| should yield \lstinline|7.50|.

  \item \textbf{Infinite Primes.}
    Write a function that returns a lazy sequence (using the \lstinline|fn| module) of prime numbers.
    Use it to print the first 20 primes.
    \emph{Hint}: use \lstinline|F.select| with a primality-testing closure over \lstinline|F.range(2, 9999999)|.

  \item \textbf{Interleave.}
    Write a function \lstinline|interleave(seq1, seq2)| that takes two lazy sequences and produces a new sequence that alternates elements: first from \lstinline|seq1|, then from \lstinline|seq2|, then \lstinline|seq1|, and so on.
    Stop when either sequence is exhausted.
\end{enumerate}

%% ───────────────────────────────────────────────────────────
\bigskip
\begin{center}\rule{0.5\linewidth}{0.4pt}\end{center}
\medskip

\noindent\textbf{What's Next}\quad
Iterators give you a powerful way to express data transformations as pipelines.
In the next chapter, we take this further: \cref{ch:functional-programming} introduces Lattice's functional programming tools---\lstinline|pipe()|, \lstinline|compose()|, currying, and the full \lstinline|fn| standard library---that let you build elegant, composable abstractions on top of the iterator foundation we have laid here.
